
\documentclass[10pt, twocolumn]{article}
\usepackage{graphicx}
\usepackage{amsmath}
\usepackage{amssymb}
\usepackage{amsfonts}
\usepackage{textcomp}
\usepackage{amsthm}

%% ---- Theorem Environments ---------------------------------
\theoremstyle{plain}
\newtheorem{theorem}{Theorem}[section]
\newtheorem{lemma}[theorem]{Lemma}
\newtheorem{corollary}[theorem]{Corollary}
\newtheorem{proposition}[theorem]{Proposition}

\theoremstyle{definition}
\newtheorem{definition}[theorem]{Definition}
\newtheorem{example}[theorem]{Example}

\theoremstyle{remark}
\newtheorem{remark}[theorem]{Remark}

\title{P-02}
\author{A. E. Mahrouss\\amlal@nekernel.org}
\date{\today}

\begin{document}
\maketitle

\section{Theorems}
\begin{theorem}
    There exists in $\mathbb{R}$
    \begin{equation}
        \frac{\int (np) \text{ } dp}{\int (px) \text{ } dx} = \frac{p^2n+2C}{px^2+2C}
    \end{equation}
    \begin{equation}
        \lim_{x\to\infty} \frac{p^2n+2C}{px^2+2C} = \frac{\operatorname{undefined}}{c} = 0
    \end{equation}
\end{theorem}

\end{document}